Table~\ref{tab:main} gives the main comparison (5 seeds, $n_{\rm train}=5000$, tuned hyperparameters) and Table~\ref{tab:cmp} the paired tests. VPT is in Lyapunov times with a cap of 15; a value is the mean over 20 test initial conditions, and tables give mean $\pm$ std over seeds.

\begin{table*}[t]\centering\caption{Main results, $n_{\rm train}=5000$, 5 seeds (mean $\pm$ std over seeds). \texttt{vpt\_min}: shortest VPT over the 20 initial conditions. ``True system'': the ODE integrated from the same warm-ups (sampling floor of the climate metrics; its VPT is the cap). No best/second marks are printed because several columns are tied.}\label{tab:main}
\small\begin{tabular}{llccccccc}
\toprule
Method & Task & vpt $\uparrow$ & vpt\_min $\uparrow$ & clim\_w1 $\downarrow$ & clim\_ok $\uparrow$ & blowup $\downarrow$ & fit (s) $\downarrow$ & $n$ \\
\midrule
True system & L63 & 15.005 $\pm$ 0.000 & 15.005 $\pm$ 0.000 & 0.048 $\pm$ 0.018 & 0.960 $\pm$ 0.055 & 0.000 $\pm$ 0.000 & 0.000 $\pm$ 0.000 & 5 \\
MLP (delay) & L63 & 5.578 $\pm$ 0.174 & 2.150 $\pm$ 0.851 & 0.063 $\pm$ 0.012 & 0.810 $\pm$ 0.139 & 0.000 $\pm$ 0.000 & 28.667 $\pm$ 1.138 & 5 \\
GRU & L63 & 3.787 $\pm$ 0.612 & 1.303 $\pm$ 0.492 & 0.058 $\pm$ 0.014 & 0.860 $\pm$ 0.108 & 0.000 $\pm$ 0.000 & 57.709 $\pm$ 2.982 & 5 \\
ESN & L63 & 9.546 $\pm$ 0.318 & 6.219 $\pm$ 1.187 & 0.061 $\pm$ 0.015 & 0.850 $\pm$ 0.141 & 0.000 $\pm$ 0.000 & 0.152 $\pm$ 0.013 & 5 \\
\midrule
True system & L96 & 14.996 $\pm$ 0.000 & 14.996 $\pm$ 0.000 & 0.036 $\pm$ 0.004 & 1.000 $\pm$ 0.000 & 0.000 $\pm$ 0.000 & 0.000 $\pm$ 0.000 & 5 \\
MLP (delay) & L96 & 2.882 $\pm$ 0.262 & 1.123 $\pm$ 0.370 & 0.036 $\pm$ 0.002 & 1.000 $\pm$ 0.000 & 0.000 $\pm$ 0.000 & 16.110 $\pm$ 0.866 & 5 \\
GRU & L96 & 1.330 $\pm$ 0.231 & 0.104 $\pm$ 0.075 & 0.045 $\pm$ 0.004 & 1.000 $\pm$ 0.000 & 0.000 $\pm$ 0.000 & 62.663 $\pm$ 2.370 & 5 \\
ESN & L96 & 2.115 $\pm$ 0.345 & 0.718 $\pm$ 0.156 & 0.040 $\pm$ 0.003 & 1.000 $\pm$ 0.000 & 0.000 $\pm$ 0.000 & 0.158 $\pm$ 0.009 & 5 \\
\bottomrule
\end{tabular}
\end{table*}

\begin{table}[t]\centering\caption{Paired comparisons over the 5 seeds of Table~\ref{tab:main} (\texttt{rh compare}). Differences are A $-$ B (positive: A is higher); $p$ is the paired $t$-test, n/a where the per-seed differences have zero variance.}\label{tab:cmp}
\small\setlength{\tabcolsep}{3pt}\begin{tabular}{lllrrrrr}\toprule
 &  &  & \multicolumn{3}{c}{vpt} & \multicolumn{2}{c}{clim\_ok}\\
Task & A & B & A $-$ B & A/B & $p$ & A $-$ B & $p$\\
\midrule
L63 & ESN & MLP & +3.97 & 1.71 & 3.6e-06 & +0.04 & 0.654\\
L63 & ESN & GRU & +5.76 & 2.52 & 3.1e-05 & -0.01 & 0.704\\
L63 & ESN & ESN no sq. & +1.41 & 1.17 & 0.008 & +0.01 & 0.799\\
L63 & GRU & GRU no noise & +0.12 & 1.03 & 0.657 & +0.23 & 0.205\\
\midrule
L96 & ESN & MLP & -0.77 & 0.73 & 0.018 & +0.00 & n/a\\
L96 & ESN & GRU & +0.78 & 1.59 & 0.002 & +0.00 & n/a\\
L96 & ESN & ESN no sq. & +1.15 & 2.20 & 0.002 & +1.00 & n/a\\
L96 & GRU & GRU no noise & +0.00 & 1.00 & 0.995 & +0.99 & 6.2e-08\\
\bottomrule\end{tabular}
\end{table}

\textbf{H1 (ESN longest VPT at 5000 samples): supported on Lorenz-63, refuted on Lorenz-96.} On Lorenz-63 the ESN's mean VPT is 1.71 times the MLP's and 2.52 times the GRU's, and both paired tests give $p<10^{-4}$ (Table~\ref{tab:cmp}). On Lorenz-96 the ESN is ahead of the GRU ($p=0.002$) but behind the MLP: the difference is $-0.77$ Lyapunov times with paired $p=0.018$. The registered rule refutes H1 on a task when a baseline's mean is at least the ESN's, so H1 is refuted as stated. Two caveats apply. First, this is the ESN at $N=500$; in the size sweep (Table~\ref{tab:rho}, seeds 0--2) the $N=2000$ reservoir reaches 3.28 on Lorenz-96, above the MLP's 2.72 on those seeds, and we did not vary the MLP's width. Second, the MLP row is not the MLP's best either: with more optimiser steps than the validation grid selected, its Lorenz-96 mean rises further (Table~\ref{tab:steps}). The ESN trains in a fraction of a second (0.152 and 0.158~s), against 16--29~s for the MLP and about one minute for the GRU (Table~\ref{tab:main}).

\textbf{H2 (VPT and climate rankings differ): met formally, not evidence of a dissociation.} On Lorenz-63 the order by mean \texttt{clim\_ok} is GRU, ESN, MLP (0.860, 0.850, 0.810) and the order by VPT is ESN, MLP, GRU, so the registered rule (orders differ on at least one task) is satisfied. We do not read this as support: the three \texttt{clim\_ok} means lie within one seed std of each other, the paired tests between them give $p=0.654$ and $p=0.704$ (Table~\ref{tab:cmp}), and on Lorenz-96 all three are tied at 1. What the data show is that at the main setting the three systems are indistinguishable in this climate metric, although their VPTs differ by a factor of up to 2.52. On Lorenz-63 all three have a \texttt{clim\_ok} somewhat below the true system's 0.960; the seed spreads overlap and we did not test that difference. The two criteria do separate within single systems (Sec.~\ref{sec:ablations}).

\textbf{Where methods fail.} The worst initial condition (\texttt{vpt\_min}, Table~\ref{tab:main}) is far below the mean for every system; the GRU's worst case on Lorenz-96 is about a tenth of a Lyapunov time. On Lorenz-96 no system reaches three Lyapunov times at this model size and data length, and the GRU is the weakest system on both tasks. No tuned system left the attractor region in a free run (\texttt{blowup}$=0$).

\textbf{Training length (H5).} Table~\ref{tab:ntrain} and Fig.~\ref{fig:ntrain} give VPT versus $n_{\rm train}$ (seeds 0--2; the MLP and GRU keep the optimiser setting tuned at 5000 samples). At $n_{\rm train}=500$ the ESN has the highest mean VPT on both tasks, so H5 is supported by the registered mean-ordering rule. The evidence is weak: the paired tests of the ESN against the MLP at 500 samples give $p=0.228$ (Lorenz-63) and $p=0.372$ (Lorenz-96), and only the GRU is clearly behind. On Lorenz-63 the ESN leads the MLP by more than 3 Lyapunov times from 2000 samples on ($p\le0.004$); its own mean changes little between 5000 and 10000 samples (within one seed std), and the MLP and GRU means show no trend from 2000 samples on. On Lorenz-96 the MLP's mean is above the ESN's from 2000 samples on, but with 3 seeds the paired test reaches 0.05 only at 10000 samples ($p=0.016$). There the MLP and GRU means rise at every step of the grid, by more than their seed std between 5000 and 10000 samples, so the data these two need on Lorenz-96 is not bracketed by this sweep; the ESN's increase over the same step is within its seed std.

\begin{table*}[t]\centering\caption{VPT (mean $\pm$ std over seeds 0--2) versus training samples. $\Delta$ = ESN $-$ MLP and $p$ its paired $t$-test.}\label{tab:ntrain}
\small\setlength{\tabcolsep}{4pt}\begin{tabular}{rcccccccccc}\toprule
$n_{\rm train}$ & L63 ESN & L63 MLP & L63 GRU & L63 $\Delta$ & L63 $p$ & L96 ESN & L96 MLP & L96 GRU & L96 $\Delta$ & L96 $p$\\
\midrule
500 & 5.92 $\pm$ 0.13 & 5.20 $\pm$ 0.84 & 2.53 $\pm$ 0.63 & +\rhval{cmp/nt500/mlp-delay/lorenz63/vpt/delta:2} & \rhval{cmp/nt500/mlp-delay/lorenz63/vpt/paired_p:3} & 0.67 $\pm$ 0.13 & 0.58 $\pm$ 0.05 & 0.05 $\pm$ 0.04 & +\rhval{cmp/nt500/mlp-delay/lorenz96/vpt/delta:2} & \rhval{cmp/nt500/mlp-delay/lorenz96/vpt/paired_p:3}\\
2000 & 9.04 $\pm$ 0.26 & 5.84 $\pm$ 0.20 & 3.68 $\pm$ 0.65 & +\rhval{cmp/nt2000/mlp-delay/lorenz63/vpt/delta:2} & \rhval{cmp/nt2000/mlp-delay/lorenz63/vpt/paired_p:3} & 1.78 $\pm$ 0.31 & 1.98 $\pm$ 0.18 & 0.42 $\pm$ 0.04 & \rhval{cmp/nt2000/mlp-delay/lorenz96/vpt/delta:2} & \rhval{cmp/nt2000/mlp-delay/lorenz96/vpt/paired_p:3}\\
5000 & 9.60 $\pm$ 0.31 & 5.55 $\pm$ 0.24 & 3.60 $\pm$ 0.72 & +\rhval{cmp/nt5000/mlp-delay/lorenz63/vpt/delta:2} & \rhval{cmp/nt5000/mlp-delay/lorenz63/vpt/paired_p:sci1} & 2.04 $\pm$ 0.46 & 2.72 $\pm$ 0.09 & 1.20 $\pm$ 0.20 & \rhval{cmp/nt5000/mlp-delay/lorenz96/vpt/delta:2} & \rhval{cmp/nt5000/mlp-delay/lorenz96/vpt/paired_p:3}\\
10000 & 9.84 $\pm$ 0.35 & 5.65 $\pm$ 0.26 & 3.58 $\pm$ 0.39 & +\rhval{cmp/nt10000/mlp-delay/lorenz63/vpt/delta:2} & \rhval{cmp/nt10000/mlp-delay/lorenz63/vpt/paired_p:3} & 2.20 $\pm$ 0.20 & 3.19 $\pm$ 0.06 & 1.74 $\pm$ 0.05 & \rhval{cmp/nt10000/mlp-delay/lorenz96/vpt/delta:2} & \rhval{cmp/nt10000/mlp-delay/lorenz96/vpt/paired_p:3}\\
\bottomrule\end{tabular}
\end{table*}

\begin{figure}[t]\centering\includegraphics[width=\columnwidth]{sweep_ntrain_fig.pdf}
\caption{VPT versus number of training samples (log axis), seeds 0--2.}\label{fig:ntrain}\end{figure}
